\begin{propositionbox}{Proposition (Left Zero Absorption)}
\begin{proposition}[Left Zero Absorption]
\label{prop:left-zero-absorption}
When multiplication distributes over addition and $0$ is the additive identity,
\[
0\cdot a=0.
\]
\end{proposition}
\end{propositionbox}

\begin{dependencies}
\begin{itemize}
  \item \hyperref[def:two-sided-distributivity]{Two-Sided Distributivity}
  \item \hyperref[def:two-sided-identity]{Two-Sided Identity}
  \item \hyperref[def:absorbing-element]{Absorbing Element}
\end{itemize}
\end{dependencies}

\begin{propositionbox}{Proposition (Right Zero Absorption)}
\begin{proposition}[Right Zero Absorption]
\label{prop:right-zero-absorption}
When multiplication distributes over addition and $0$ is the additive identity,
\[
a\cdot 0=0.
\]
\end{proposition}
\end{propositionbox}

\begin{dependencies}
\begin{itemize}
  \item \hyperref[def:two-sided-distributivity]{Two-Sided Distributivity}
  \item \hyperref[def:two-sided-identity]{Two-Sided Identity}
  \item \hyperref[def:absorbing-element]{Absorbing Element}
\end{itemize}
\end{dependencies}

\begin{propositionbox}{Proposition (Zero Absorption)}
\begin{proposition}[Zero Absorption]
\label{prop:zero-absorption}
When multiplication distributes over addition and $0$ is the additive identity,
\[
0\cdot a=0
\qquad\text{and}\qquad
a\cdot 0=0.
\]
\end{proposition}
\end{propositionbox}

\begin{dependencies}
\begin{itemize}
  \item \hyperref[prop:left-zero-absorption]{Left Zero Absorption}
  \item \hyperref[prop:right-zero-absorption]{Right Zero Absorption}
\end{itemize}
\end{dependencies}

\begin{propositionbox}{Proposition (Left Sign Rule)}
\begin{proposition}[Left Sign Rule]
\label{prop:left-sign-rule}
When multiplication distributes over addition and additive inverses exist,
\[
(-a)\cdot b=-(a\cdot b).
\]
\end{proposition}
\end{propositionbox}

\begin{dependencies}
\begin{itemize}
  \item \hyperref[def:two-sided-distributivity]{Two-Sided Distributivity}
  \item \hyperref[def:two-sided-inverse]{Two-Sided Inverse}
  \item \hyperref[prop:left-zero-absorption]{Left Zero Absorption}
\end{itemize}
\end{dependencies}

\begin{propositionbox}{Proposition (Right Sign Rule)}
\begin{proposition}[Right Sign Rule]
\label{prop:right-sign-rule}
When multiplication distributes over addition and additive inverses exist,
\[
a\cdot(-b)=-(a\cdot b).
\]
\end{proposition}
\end{propositionbox}

\begin{dependencies}
\begin{itemize}
  \item \hyperref[def:two-sided-distributivity]{Two-Sided Distributivity}
  \item \hyperref[def:two-sided-inverse]{Two-Sided Inverse}
  \item \hyperref[prop:right-zero-absorption]{Right Zero Absorption}
\end{itemize}
\end{dependencies}

\begin{propositionbox}{Proposition (Sign Rule)}
\begin{proposition}[Sign Rule]
\label{prop:sign-rule}
When multiplication distributes over addition and additive inverses exist,
\[
(-a)\cdot b=-(a\cdot b)
\qquad\text{and}\qquad
a\cdot(-b)=-(a\cdot b).
\]
\end{proposition}
\end{propositionbox}

\begin{dependencies}
\begin{itemize}
  \item \hyperref[prop:left-sign-rule]{Left Sign Rule}
  \item \hyperref[prop:right-sign-rule]{Right Sign Rule}
\end{itemize}
\end{dependencies}

\begin{corollarybox}{Corollary (Negative One Rule)}
\begin{corollary}[Negative One Rule]
\label{cor:negative-one-rule}
When multiplication distributes over addition and additive inverses exist,
\[
(-1)\cdot a=-a.
\]
\end{corollary}
\end{corollarybox}

\begin{dependencies}
\begin{itemize}
  \item \hyperref[prop:left-sign-rule]{Left Sign Rule}
\end{itemize}
\end{dependencies}

\begin{propositionbox}{Proposition (Product of Negatives)}
\begin{proposition}[Product of Negatives]
\label{prop:product-of-negatives}
When multiplication distributes over addition and additive inverses exist,
\[
(-a)\cdot(-b)=a\cdot b.
\]
\end{proposition}
\end{propositionbox}

\begin{dependencies}
\begin{itemize}
  \item \hyperref[prop:left-sign-rule]{Left Sign Rule}
  \item \hyperref[prop:right-sign-rule]{Right Sign Rule}
\end{itemize}
\end{dependencies}

\begin{definitionbox}{Definition (Zero Divisor)}
\begin{definition}[Zero Divisor]
\label{def:zero-divisor}
Let $K$ be an arithmetic system with multiplication and zero. A nonzero
element $a\in K$ is a \emph{zero divisor} if there exists a nonzero element
$b\in K$ such that $ab=0$.
\end{definition}
\end{definitionbox}

\begin{dependencies}
\begin{itemize}
  \item \hyperref[def:binary-arithmetic-operation]{Binary Arithmetic Operation}
\end{itemize}
\end{dependencies}

\begin{theorembox}{Theorem (No Zero Divisors)}
\begin{theorem}[No Zero Divisors]
\label{thm:no-zero-divisors}
In an arithmetic system with no nonzero zero divisors,
\[
a\cdot b=0 \Longrightarrow a=0 \lor b=0.
\]
\end{theorem}
\end{theorembox}

\begin{dependencies}
\begin{itemize}
  \item \hyperref[def:binary-arithmetic-operation]{Binary Arithmetic Operation}
\end{itemize}
\end{dependencies}
